% Definition file for row def_3_15 -- "CollidersAndNon".
% Auto-generated stub by scaffold/solve_chapter.py at row start. The
% manager agent (or a worker dispatched by it) fills the body below with the
% Lean-faithful definition statement(s). Stays close to the lecture notes.
%
% This file is a sibling of `main.tex` inside `<subsection>/tex/`. The
% `subfiles` package lets it compile both standalone (resolving its
% preamble via `main.tex`) and as part of `main.tex` via `\subfile{...}`.

\documentclass[main]{subfiles}

\begin{document}
% ref: def_3_15
% title: CollidersAndNon
\def\rowref{def\_3\_15}
\phantomsection\label{def_3_15}

\begin{Def}[Colliders and non-colliders]\label{def:collider_noncollider}
    Let $G = (J, V, E, L)$ be a CDMG (def \ref{def-cdmg}); throughout we use the notation of def \ref{not-cdmg} and the edge classifications of def \ref{def-edge-relations}. Let
    \[
        \pi = (v_0, a_0, v_1, a_1, \dots, v_{n-1}, a_{n-1}, v_n)
    \]
    be a walk in $G$ (def \ref{def:walks}, item~i.) for some integer $n \ge 0$, with $v_0, \dots, v_n \in J \cup V$ and $a_0, \dots, a_{n-1} \in E \cup L$ satisfying the walk constraint at every $k \in \{0, 1, \dots, n - 1\}$:
    \[
        \bigl(a_k = (v_k, v_{k+1}) \;\land\; a_k \in E \cup L\bigr)
        \;\lor\;
        \bigl(a_k = (v_{k+1}, v_k) \;\land\; a_k \in E\bigr).
    \]

    \emph{Arrowhead count at a position.}
    For a position $k \in \{0, 1, \dots, n\}$, the \emph{arrowhead count at $v_k$ on $\pi$}, denoted $\mathrm{ah}_\pi(k)$, is the cardinality
    \[
        \mathrm{ah}_\pi(k)
            := \Bigl|\Bigl\{\, i \in \{k - 1, k\} \,\Bigm|\, 0 \le i \le n - 1 \,\land\, a_i \text{ is an edge into } v_k \text{ in } G \,\Bigr\}\Bigr|,
    \]
    where ``$a_i$ is an edge into $v_k$ in $G$'' is in the sense of def \ref{def-edge-relations}, item~ii.; spelled out, for any $i \in \{0, 1, \dots, n - 1\}$ writing $a_i = (e_1, e_2)$,
    \[
        a_i \text{ is an edge into } v_k
        \quad\iff\quad
        \bigl(a_i \in E \;\land\; e_2 = v_k\bigr)
            \;\lor\;
        \bigl(a_i \in L \;\land\; (e_1 = v_k \,\lor\, e_2 = v_k)\bigr).
    \]
    The index set defining $\mathrm{ah}_\pi(k)$ ranges only over the (at most two) walk-incident edges at position $k$, namely $a_{k-1}$ (when $k \ge 1$) and $a_k$ (when $k \le n - 1$); edges of $G$ that are not incident to $v_k$ along $\pi$ are excluded by construction. At an end-position $k \in \{0, n\}$ at most one such index is admissible, so $\mathrm{ah}_\pi(k) \le 1$ automatically (and $\mathrm{ah}_\pi(0) = 0$ on the trivial walk $n = 0$, where neither index is admissible). At an interior position $1 \le k \le n - 1$ both indices are admissible, so $\mathrm{ah}_\pi(k) \in \{0, 1, 2\}$. (The stored-pair test ``$a_i$ is an edge into $v_k$'' above is ambiguous at a walk-incident directed self-loop $a_i \in E$ with $v_i = v_{i+1}$, where its $E$-clause fires unconditionally at both adjacent positions $k \in \{i, i + 1\}$; the canonical disambiguation at such positions is the side-aware reading of the addition tag \texttt{[collider\_side\_aware\_walkstep\_predicates]}, spelled out in the Treatment of directed self-loops paragraph below.)

    \emph{Encoding note: side-aware reading via a typed walk-step representation.}
    At a walk-incident edge $a_i$ with distinct endpoints ($v_i \ne v_{i+1}$), the two walk-constraint disjuncts of def \ref{def:walks}, item~i.\ are mutually exclusive, so whichever disjunct fires unambiguously determines the walk-traversal direction at $a_i$ (forward $v_i \to v_{i+1}$ for the first disjunct, backward $v_{i+1} \to v_i$ for the second). Note that def \ref{def-cdmg} explicitly admits \emph{writing-mirror} pairs: the irreflexivity restriction applies to $L$ only, and no graph-theoretic mutual-exclusion is imposed between the directed channel $E$ and the bidirected channel $L$, so a single underlying vertex pair $\{u, v\}$ may simultaneously support a directed edge $(u, v) \in E$ (or $(v, u) \in E$) \emph{and} a bidirected edge $s(u, v) \in L$. On a non-self-loop walk-step $a_i$ the side-aware reading of the addition tag \texttt{[collider\_side\_aware\_walkstep\_predicates]} therefore reads, at each adjacent position $k \in \{i, i+1\}$, the \emph{disjunction over channels} of (a) the constructor-tagged channel's directional contribution -- $1$ at the walk-traversal target and $0$ at the source for a directed $a_i \in E$; mirror for a directed step taken backward; $1$ at both adjacent positions for a bidirected $a_i \in L$ -- and (b) an additional $L$-channel head-contribution at the otherwise-zero side whenever the same stored pair is also in $L$ (the writing-mirror case). With this OR-of-channels convention, the side-aware reading coincides pointwise with the literal stored-pair ``$a_i$ is an edge into $v_k$'' test of def \ref{def-edge-relations}, item~ii.\ at every non-self-loop walk-step: a writing-mirror pair contributes from both channels under both readings, a pure-$E$ pair contributes from the $E$-clause under both, and a pure-$L$ pair contributes from the $L$-clause under both. At a walk-incident self-loop edge $a_i \in E$ with $v_i = v_{i+1} = v$, by contrast, both walk-constraint disjuncts apply syntactically and the stored-pair test is ambiguous -- firing at both adjacent positions, as discussed in the Treatment of directed self-loops paragraph below -- and the side-aware reading uses the walk-step's explicitly-recorded direction tag (forward or backward) to disambiguate, taking precedence over the stored-pair test per the addition; the $L$-channel disjunct cannot fire at a self-loop because def \ref{def-cdmg}'s irreflexivity restriction on $L$ rules out $s(v, v) \in L$ outright, so the side-aware refinement at self-loops is preserved by construction.

    The Lean realisation of this definition implements the side-aware reading faithfully via a typed walk-step representation: each walk-step carries a constructor recording whether the underlying $E$-edge is traversed forward or backward, with a separate constructor for the bidirected $L$-case. The per-edge contribution to $\mathrm{ah}_\pi(k)$ is then read off the constructor tag and the walk-step's source/target type indices directly, without consulting any node-equality between $v_k$ and the ordered-pair components; an additional opposite-side $L$-channel head-contribution disjunct $s(u, v) \in L$ on the constructor-tagged directed branches re-fires the $L$-channel reading at writing-mirror pairs (vacuous at self-loops by the irreflexivity restriction on $L$). This yields the side-aware OR-of-channels contribution uniformly across self-loop and non-self-loop walk-steps.

    \emph{Classification.}
    The position $k \in \{0, 1, \dots, n\}$ on the walk $\pi$ -- or, by abuse, the node $v_k$ at position $k$ -- is called:
    \begin{enumerate}[label=\roman*.)]
        \item a \emph{non-collider on $\pi$} iff $\mathrm{ah}_\pi(k) \le 1$;
        \item a \emph{collider on $\pi$} iff $\mathrm{ah}_\pi(k) = 2$.
    \end{enumerate}
    By construction these two classifications are mutually exclusive and jointly exhaustive over $\{0, 1, \dots, n\}$: every position $k$ on $\pi$ is exactly one of a non-collider or a collider on $\pi$.

    \emph{Reconciliation with the source-block pattern writings.}
    The source block characterises non-colliders by a list of four pattern-cases (end-node, left chain, right chain, fork) glossed as ``at most one arrowhead pointing towards $v_k$'', and colliders by a single pattern-case glossed as ``two arrowheads pointing towards $v_k$ on the walk $\pi$''. Below we reconcile these visual writings with the count-based classification, taking the count-based classification (the LN's ``i.e.''\ gloss) as canonical. The visual writings $v_{k-1} \tuh v_k$, $v_{k-1} \hut v_k$, $v_{k-1} \huh v_k$, $v_{k-1} \suh v_k$, $v_{k-1} \hus v_k$ are unfolded per def \ref{not-cdmg}, items~2--6.
    \begin{itemize}
        \item \textbf{end-node} ($k \in \{0, n\}$): $\mathrm{ah}_\pi(k) \le 1$ holds automatically, hence the position is a non-collider on $\pi$. (For the trivial walk $n = 0$, the unique position $k = 0 = n$ has $\mathrm{ah}_\pi(0) = 0$ and is a non-collider.)
        \item \textbf{left chain} ($1 \le k \le n - 1$; LN writing $v_{k-1} \hut v_k \hus v_{k+1}$): the walk's $a_{k-1}$ is \emph{not} an edge into $v_k$, while its $a_k$ \emph{is} an edge into $v_k$. Spelled out via the walk constraint, $a_{k-1} = (v_k, v_{k-1}) \in E$ (a directed edge with arrowhead at $v_{k-1}$ and tail at $v_k$), and either $a_k = (v_{k+1}, v_k) \in E$ or $a_k = (v_k, v_{k+1}) \in L$. Count: $0 + 1 = 1$; non-collider.
        \item \textbf{right chain} ($1 \le k \le n - 1$; LN writing $v_{k-1} \suh v_k \tuh v_{k+1}$): the walk's $a_{k-1}$ \emph{is} an edge into $v_k$, while its $a_k$ is \emph{not} an edge into $v_k$. Spelled out, either $a_{k-1} = (v_{k-1}, v_k) \in E$ or $a_{k-1} = (v_{k-1}, v_k) \in L$, and $a_k = (v_k, v_{k+1}) \in E$. Count: $1 + 0 = 1$; non-collider.
        \item \textbf{fork} ($1 \le k \le n - 1$; LN writing $v_{k-1} \hut v_k \tuh v_{k+1}$): neither walk-incident edge is an edge into $v_k$. Spelled out, $a_{k-1} = (v_k, v_{k-1}) \in E$ and $a_k = (v_k, v_{k+1}) \in E$. Count: $0 + 0 = 0$; non-collider.
        \item \textbf{collider} ($1 \le k \le n - 1$; LN writing $v_{k-1} \suh v_k \hus v_{k+1}$): both walk-incident edges $a_{k-1}, a_k$ are edges into $v_k$. Spelled out, $a_{k-1} \in E \cup L$ with right-component $v_k$ if directed (or either component $v_k$ if bidirected), and analogously for $a_k$ with left-component $v_k$ if directed (or either component $v_k$ if bidirected). Count: $1 + 1 = 2$; collider.
    \end{itemize}
    The four non-collider patterns realise exactly the interior counts $\{0, 1\}$, and the collider pattern realises exactly the interior count $2$; together with the end-node case they exhaust the positions of $\pi$. The pattern-cases consult \emph{walk-incident} edges $a_{k-1}, a_k$ only -- arrowheads contributed by edges of $G$ that are not on $\pi$ at $v_k$ are excluded from $\mathrm{ah}_\pi(k)$ by construction. (This fixes the scope of the LN's gloss ``at most one (resp.\ two) arrowheads pointing towards $v_k$''.) For walks that traverse a directed self-loop $a_i \in E$ with $v_i = v_{i+1} = v$ -- where the literal pattern writings would otherwise admit any of the four conjunctions above simultaneously -- the recorded walk-traversal direction of $a_i$ (forward or backward) selects exactly one pattern per adjacent position, per the side-aware reading of the addition tag \texttt{[collider\_side\_aware\_walkstep\_predicates]}; see the Treatment of directed self-loops paragraph below for the per-position contribution at a self-loop walk-step.

    \emph{Treatment of directed self-loops.}
    A directed self-loop $(v, v) \in E$ is admitted by def \ref{def-cdmg} (the irreflexivity restriction applies to $L$, not to $E$). At a walk-incident self-loop edge $a_i = (v_i, v_{i+1}) = (v, v) \in E$ with $v_i = v_{i+1} = v$, both the LN's literal pattern writings $v_{k-1} \tuh v_k$ and $v_{k-1} \hut v_k$ (def \ref{not-cdmg}, items~2--3, unfolding to $(v_{k-1}, v_k) \in E$ and $(v_k, v_{k-1}) \in E$ respectively) and the stored-pair ``$a_i$ is an edge into $v_k$'' test of def \ref{def-edge-relations}, item~ii.\ are syntactically ambiguous: they collapse to the single condition $(v, v) \in E$, which fires unconditionally at both adjacent positions $k \in \{i, i + 1\}$ for which $v_k = v$. In particular, the stored-pair count formula above would commit, at such a self-loop traversal, to contribution $1$ from $a_i$ at both adjacent positions, and a single interior position $k$ traversing the self-loop would simultaneously match the ``left chain'', ``right chain'', ``fork'', and ``collider'' pattern conjunctions of the Reconciliation paragraph above.

    The canonical disambiguation adopted in this definition is the \emph{side-aware} reading introduced by the addition tag \texttt{[collider\_side\_aware\_walkstep\_predicates]}: each walk-step $a_i$ on $\pi$ carries a walk-traversal direction (forward $v_i \to v_{i+1}$ when $a_i$ satisfies the first disjunct of the walk constraint of def \ref{def:walks}, item~i.; backward $v_{i+1} \to v_i$ when it satisfies the second), with the choice recorded explicitly at a directed self-loop where both disjuncts apply syntactically. A directed self-loop $a_i \in E$ then contributes $1$ to $\mathrm{ah}_\pi(k)$ at its walk-traversal target position only -- $k = i + 1$ when the recorded direction is forward, $k = i$ when backward -- and $0$ at its walk-traversal source position. This explicitly replaces, at self-loops, the prior count-based commitment to contribution $1$ at both adjacent positions; on walks that traverse no directed self-loops the side-aware and stored-pair readings coincide, so no other part of the definition is affected. The recorded walk-traversal direction also selects exactly one of the four literal $\tuh / \hut$ pattern conjunctions per self-loop walk-step at each adjacent position, resolving the apparent multi-pattern overlap noted above.
\end{Def}

\end{document}
